\chapter{Introducao}
\label{cap:introducao}

A introducao apresenta o tema, delimita o problema e anuncia o que vem pela
frente. Escreva-a por ultimo: so depois de terminar o trabalho voce sabe o que
esta introduzindo.

\section{Contextualizacao}
\label{sec:contextualizacao}

Situe o leitor no assunto, partindo do geral para o especifico. Toda afirmacao
que nao seja sua precisa vir com a fonte, seja entre parenteses \cite{nbr14724},
seja no corpo da frase, como em \citeonline{abntex2}.

\section{Problema de pesquisa}
\label{sec:problema}

Enuncie o problema como uma pergunta clara e respondivel.

\section{Objetivos}
\label{sec:objetivos}

\subsection{Objetivo geral}

O que o trabalho se propoe a alcancar, em uma frase iniciada por um verbo no
infinitivo.

\subsection{Objetivos especificos}

\begin{alineas}
  \item primeiro passo para chegar ao objetivo geral;
  \item segundo passo;
  \item terceiro passo.
\end{alineas}

\section{Justificativa}
\label{sec:justificativa}

Por que este trabalho precisa existir: a lacuna que ele preenche e a quem
interessa preenche-la.

\section{Estrutura do trabalho}
\label{sec:estrutura}

O Capitulo~\ref{cap:referencial} apresenta o referencial teorico. O
Capitulo~\ref{cap:metodologia} descreve a metodologia. O
Capitulo~\ref{cap:resultados} discute os resultados, e o
Capitulo~\ref{cap:conclusao} traz as consideracoes finais.
